\section{Conclusions}

This study developed and evaluated a~voxel-sampler-based processing pipeline for RandLA-Net, using overlapping spatial windows and score interpolation to reconstruct point-wise labels. The comparison with an adapted probability-sampler-based workflow shows that input construction and reconstruction affect both segmentation quality and computational cost without changing the network architecture.

The proposed pipeline more than halved mean elapsed processing time in the DALES benchmarks, at the cost of higher memory demand. Segmentation quality depended on the dataset and class: probability sampling achieved higher overall mIoU on DALES, while the proposed workflow achieved higher mIoU on BRIK and higher F1 for buildings on DALES and for poles and support devices on BRIK. These results support its use where shorter processing time and the recovery of these classes are priorities. The density-reduction experiment further showed a~smaller decline in accuracy and complete returned-prediction coverage for the proposed workflow across the tested settings. This trend suggests that fixed spatial windows are a~promising choice under density variation, with accuracy comparisons conditioned on the points receiving predictions.

The remaining development priorities are to reduce preparation and reconstruction costs and to improve the treatment of sparse structures and window boundaries. Joint evaluation of window geometry, sampling thresholds, and score interpolation could help retain the observed speed and class-level gains while improving recall for power lines and railway tracks. Such evaluation should combine per-class quality, prediction coverage, elapsed time, and memory use to guide the choice of pipeline for the intended application.
